\clearpage
\section{Repeated Scalar Execution}\label{page:repeated-scalar-execution}

\texttt{REP Rn, (<instruction>)} and \texttt{REPcc Rn, (<instruction>)} repeat one scalar instruction. The body is
decoded and checked for repeat eligibility before a zero-count annul is applied, then remains fixed for that repeat
context. A body may not be a repeat operation, a repeat-control instruction, or a control transfer.

Entering repeat state captures the selected Rn value as an unsigned remaining count. A valid zero count performs no
body side effects. Otherwise each iteration evaluates operands and commits the body\textquotesingle{}s register, memory, and auto-update
effects according to its (:term:base.terms.ordinary:) scalar rules. A body read of the selected counter register observes the
iteration-start architectural value; a body write to that register is (:term:base.terms.ignored:).

Each successfully committed iteration decrements the captured count before an event may be admitted at the next
iteration boundary. A condition-false (:ref:base.instructions.REPcc:) iteration still commits its body and decrements the count before
repeat state is cleared. If the body faults, that iteration commits nothing and does not decrement the count; earlier
iterations remain committed. Event entry then saves the repeat-prefix address in (:ref:base.events.frame.EVENT_FRAME.SAVED_PC:) and clears hidden repeat
state, so ERET returns to the prefix and re-executes it normally with the unchanged counter.

The \texttt{REPcc observation} attribute identifies the value tested after the body executes. A
\texttt{result:operand} or \texttt{source:operand} observation names an operand; \texttt{computed} names the derived
value defined by that instruction. The temporary condition image sets Z exactly when the observed value is zero,
sets N from its most significant bit at the observation width, and clears C and V. It does not sample or overwrite
architectural FLAGS, although the body\textquotesingle{}s own FLAGS effects commit normally. \texttt{REP} is the unconditional T
spelling; condition F is (:term:base.terms.reserved:).

While repeat state is active, architectural PC identifies the body instruction. After a successful body commit, the
counter decrement completes before an event may be admitted between iterations. Both that event path and the
fault path above save the repeat-prefix address in (:ref:base.events.frame.EVENT_FRAME.SAVED_PC:) and clear hidden repeat state. ERET restores the
(:term:base.terms.ordinary:) saved (:term:base.terms.architectural_state:) and begins execution at that prefix address; it does not restore hidden
repeat (:term:base.terms.continuation:) state. The prefix is therefore decoded and executed normally, using the already-decremented
architectural Rn value. (:ref:base.registers.STATE.STATUS.TF:) and (:ref:base.registers.STATE.STATUS.RF:) apply once to each committed body iteration; a zero-count annul is also one trace
unit.
